% ---
% id: foundations/rem:summation-order
% kind: rem
% title: "Dependence of a computed sum on its order"
% summary: >-
%   Two orderings of the same three binary64 numbers have the same exact sum
%   but different computed sums, and only one of them rounds nowhere. A
%   computed sum is therefore a determinate number only once its order is fixed.
% prominence: standard
% status: accepted
% uses: [def:binary64, def:declared-order, lem:binary64-lattice, lem:exact-accumulation]
% assumes: []
% refs: []
% provenance: newly-derived
% lean:
% ---

\begin{remark}\label{rem:summation-order}
Binary64 addition is not associative, so the computed sequential sum of
Definition~\ref{def:declared-order} depends on the order of its terms. By
Lemma~\ref{lem:binary64-lattice}\,(ii) the binary64 numbers in $[1,2)$ are the
multiples of $2^{-52}$, so $1+2^{-60}$ lies strictly between the adjacent
binary64 numbers $1$ and $1+2^{-52}$ and is nearer to $1$; hence
$\mathrm{fl}(1+2^{-60})=1$ (Definition~\ref{def:binary64}\,(ii)). The
sequence $(1,-1,2^{-60})$ therefore has computed sequential sum
$0,1,0,2^{-60}$, every addition being exact, while $(1,2^{-60},-1)$ has
computed sequential sum $0,1,1,0$, its second addition rounding. The exact sum is $2^{-60}$ in both orders. Only the first
accumulation rounds nowhere, and its computed sum is the exact sum, as
Lemma~\ref{lem:exact-accumulation} guarantees; the second computed sum is
$0$.
\end{remark}
